\chapter{Presupuesto}
\label{chap:presupuesto}

La guía pide estimar recursos, expresar costos en soles (u otra moneda acordada) e indicar \textbf{fuente de financiamiento}. En este trabajo el gasto relevante combina \textbf{materiales de cierre} (impresión, trámites), \textbf{movilidad local}, \textbf{conectividad} si hiciera falta para pruebas remotas, y---solo si el diseño del MCP lo requiere---\textbf{APIs o nube} de bajo monto. La computadora y el software principal (Node, Ionic, editor, Jest/Jasmine) se consideran \textbf{recursos ya disponibles}; si en algún caso se exigiera licencia paga, esa partida se añade con cotización.

Los montos de la Tabla~\ref{tabla:presupuesto} están \textbf{afinados como orden de magnitud} para un plan desarrollado en Perú en 2026 (impresión, movilidad en ciudad, pequeños gastos de servicio). Deben \textbf{validarse y ajustarse} con la asesora y, de ser necesario, con la secretaría de la escuela antes de aprobar el plan definitivo.

\begin{table}[h]
\centering
\small
\caption{Presupuesto referencial del proyecto de tesis (S/.), a validar con asesoría.}
\label{tabla:presupuesto}
\resizebox{\linewidth}{!}{%
\begin{tabular}{p{5cm}r p{3.6cm}}
\toprule
\textbf{Componente} & \textbf{Costo (S/.)} & \textbf{Fuente} \\
\midrule
Impresión del plan de tesis, borradores y versión final de tesis; encuadernación & 280--450 & Personal \\
Trámites y pagos administrativos (inscripción, constancias, otros según facultad) & 120--280 & Personal \\
Movilidad local (asesorías en campus, biblioteca, sustentación) & 180--350 & Personal \\
Conectividad reforzada durante picos de desarrollo (opcional; ej.\ datos móviles) & 0--120 & Personal \\
APIs de modelo o tokens (solo si la capa MCP los usa en experimentos acotados; muchos escenarios pueden quedar en 0) & 0--180 & Personal \\
Servicios cloud opcionales (p.\,ej.\ CI o análisis estático de pago; SonarCloud u homólogos suelen tener capa gratuita) & 0--150 & Personal \\
Contingencia (materiales no previstos, ajustes de última hora) & 120--200 & Personal \\
\midrule
\textbf{Total referencial} & \multicolumn{2}{r}{\textbf{700--1\,730}} \\
\bottomrule
\end{tabular}%
}
\end{table}

Si se confirma con la asesora que \textbf{no} habrá costo por APIs ni servicios cloud, el total tiende a situarse entre \textbf{700 y 1\,200}~S/. aproximadamente (impresión, trámites, movilidad y contingencia). El rango superior reserva margen para cotizaciones reales de copistería y tasas vigentes al momento de la presentación.
